\documentclass[a4paper,11pt]{article}
\usepackage{exam-style}

\fancyhead[L]{\fontsize{9pt}{10pt}\selectfont\color{ctp-subtext0}341 农业综合知识三（信电）· 模拟卷五}
\fancyhead[R]{\fontsize{9pt}{10pt}\selectfont\color{ctp-subtext0}信电学院部分}

\begin{document}
\examtitle

% ============================================================
\parttitle{第一部分 \quad 程序设计基础知识（共75分）}

\qtypename{一、填空题}{共5题，每题2分，共10分}
\anslines{0.1cm}

\q C 语言中，若 \code{int a = -10;} 则表达式 \code{a > 0 ? a : -a} 的值是 \fillblank 。

\q 若有 \code{int a[5] = \{1, 2, 3, 4, 5\}, *p = a + 2;} 则 \code{p[-1]} 的值是 \fillblank 。

\q 在 C 语言中，函数直接或间接调用自身的现象称为 \fillblank 。

\q 若定义 \code{typedef struct \{int x; int y;\} Point;} 则创建变量 \code{p} 的方式是：\fillblank 。

\q C 语言中，以只读方式打开一个已存在的文本文件，\code{fopen} 的模式字符串是 \fillblank 。

\qtypename{二、简答题}{共5题，每题3分，共15分}

\q 简述关键字 \code{const} 在 C 语言中的三种常见用法（修饰变量、修饰指针、修饰函数参数）。
\anslines{1.5cm}

\q 解释 \code{int *p[5]} 和 \code{int (*p)[5]} 的含义有何不同？
\anslines{1.5cm}

\q 简述动态内存分配函数 \code{malloc}、\code{calloc} 和 \code{realloc} 的区别。
\anslines{1.5cm}

\q 简述 \code{return} 语句在函数中的两种作用。
\anslines{1.5cm}

\q 简要解释函数 \code{fread} 和 \code{fwrite} 的参数含义和使用场景。
\anslines{1.5cm}

\qtypename{三、分析题}{共10题，每题2分，共20分}

\q
\begin{lstlisting}[language={[custom]C}]
int x = 0x1A;
printf("%d", x & 0x0F);
\end{lstlisting}
运行结果：\underline{\hspace{3cm}}


\q
\begin{lstlisting}[language={[custom]C}]
int a = 1, b = 2, c = 3;
if (a > b)
    if (a > c)
        printf("a");
else
    printf("b");
printf("c");
\end{lstlisting}
运行结果：\underline{\hspace{3cm}}


\q
\begin{lstlisting}[language={[custom]C}]
int s = 0, i = 1;
while (i <= 5) {
    s += i;
    i++;
    if (s > 6) break;
}
printf("s=%d i=%d", s, i);
\end{lstlisting}
运行结果：\underline{\hspace{3cm}}


\q
\begin{lstlisting}[language={[custom]C}]
int a[4] = {1, 1};
for (int i = 2; i < 4; i++)
    a[i] = a[i-1] + a[i-2];
for (int i = 0; i < 4; i++)
    printf("%d ", a[i]);
\end{lstlisting}
运行结果：\underline{\hspace{3cm}}


\q
\begin{lstlisting}[language={[custom]C}]
char s[] = "Hello";
char *p = s;
while (*p) p++;
printf("%ld", p - s);
\end{lstlisting}
运行结果：\underline{\hspace{3cm}}


\q
\begin{lstlisting}[language={[custom]C}]
int a = 5;
int *p = &a;
int **pp = &p;
**pp = 20;
printf("%d", a);
\end{lstlisting}
运行结果：\underline{\hspace{3cm}}


\q
\begin{lstlisting}[language={[custom]C}]
void func(int *a, int *b) {
    int *t = a; a = b; b = t;
}
int main() {
    int x = 3, y = 4;
    int *p = &x, *q = &y;
    func(p, q);
    printf("%d %d", *p, *q);
    return 0;
}
\end{lstlisting}
运行结果：\underline{\hspace{3cm}}


\q
\begin{lstlisting}[language={[custom]C}]
#define ADD(x, y) ((x) + (y))
int a = 3, b = 4;
printf("%d", ADD(a, b) * 2);
\end{lstlisting}
运行结果：\underline{\hspace{3cm}}


\q
\begin{lstlisting}[language={[custom]C}]
int i = 0, s = 0;
do {
    s += i;
    i++;
} while (i < 0);
printf("%d", s);
\end{lstlisting}
运行结果：\underline{\hspace{3cm}}


\q
\begin{lstlisting}[language={[custom]C}]
int arr[3] = {1, 2, 3};
int *p = arr;
printf("%d ", *(p+1));
printf("%d ", p[2]);
\end{lstlisting}
运行结果：\underline{\hspace{3cm}}


\qtypename{四、程序设计题}{共2题，每题15分，共30分}

\pq 编写一个 C 语言程序，实现以下功能：定义一个函数 \code{void invert(char *s)}，将字符串 s 原地逆置（例如 "hello" 逆置为 "olleh"）。然后在主函数中从键盘输入一行字符，调用逆置函数处理并输出结果。要求不能使用 \code{string.h} 中的字符串函数。
\anslines{5cm}

\pq 编写一个 C 程序，定义一个 \code{struct Date \{int year; int month; int day;\}} 结构体类型。编写函数 \code{int isLeapYear(int year)} 判断是否为闰年，\code{int daysInMonth(int year, int month)} 返回指定月份的天数，\code{int dayOfYear(struct Date d)} 计算该日期是该年的第几天。在主函数中输入一个日期，调用函数输出该日期为该年的第几天。
\anslines{6cm}

% ============================================================
\parttitle{第二部分 \quad 数据库技术与应用（共75分）}

\qtypename{一、填空题}{共5题，每题2分，共10分}

\q 数据库系统产生的主要操作异常包括插入异常、删除异常和 \fillblank 。

\q 若关系 R 和 S 的属性数分别为 4 和 3，则 R 与 S 的笛卡尔积的属性数为 \fillblank 。

\q SQL 语句中 \code{HAVING} 子句用于对 \fillblank 进行筛选。

\q 两段锁协议（2PL）分为扩展阶段和 \fillblank 阶段。

\q 在 MySQL 中，\code{SHOW DATABASES} 语句的功能是 \fillblank 。

\qtypename{二、简答题}{共2题，每题5分，共10分}

\q 解释 BCNF 与 3NF 的关系与区别。如果一个关系模式属于 BCNF，它一定属于 3NF 吗？反之呢？举例说明。
\anslines{3cm}

\q 简述数据库并发控制中加锁机制的基本思想，并说明死锁的产生原因及常见的解决策略。
\anslines{3cm}

\qtypename{三、操作题}{共8小题，共15分}

设有关系模式：\\
部门表 D(\underline{Dno}, Dname, Dlocation)\\
员工表 E(\underline{Eno}, Ename, Esex, Salary, Dno)\\
项目表 P(\underline{Pno}, Pname, Budget, Dno)\\
参与表 EP(\underline{Eno, Pno}, Hours)

\q （2分）用 SQL 查询工资高于 8000 的员工姓名及其所在部门名称。
\anslines{0.8cm}

\q （2分）用 SQL 统计每个部门中男女员工的人数。
\anslines{0.8cm}

\q （2分）用 SQL 查询预算超过 100000 且属于 "技术部" 的项目名称。
\anslines{0.8cm}

\q （2分）用关系代数表达：查询参与了所有项目的员工编号。
\anslines{1.2cm}

\q （2分）用 SQL 查询总工时（Hours）超过 100 小时的员工编号及总工时。
\anslines{0.8cm}

\q （2分）用 SQL 查询没有参与任何项目的员工姓名。
\anslines{0.8cm}

\q （1分）将员工 "E010" 的工资增加 15\%。
\anslines{0.8cm}

\q （2分）创建索引 \code{idx_dno_salary}，提高按部门编号和工资查询的效率。
\anslines{0.8cm}

\qtypename{四、分析题}{共1题，共10分}

\q 设有关系模式 R(A, B, C, D, E, F)，函数依赖集 F = \{AB $\rightarrow$ C, C $\rightarrow$ D, D $\rightarrow$ E, E $\rightarrow$ F, B $\rightarrow$ F\}。
\begin{enumerate}[leftmargin=2em, itemsep=3pt]
  \item[（1）] （4分）利用 Armstrong 公理推导出 R 的所有候选键。
  \item[（2）] （3分）该关系模式是否存在传递函数依赖？若存在，请指出。
  \item[（3）] （3分）将 R 分解为 3NF，要求保持函数依赖。
\end{enumerate}
\anslines{4cm}

\qtypename{五、设计题}{共1题，共10分}

\q 某快递公司需要设计一个快递管理系统，已知：
\begin{itemize}
  \item 每位客户有客户编号、姓名、手机号、地址；
  \item 每个快递网点有网点编号、网点名称、网点地址、联系电话；
  \item 每位快递员有工号、姓名、电话、所属网点；
  \item 每件快递有运单号、收件人信息（姓名、手机、地址）、寄件人信息（姓名、手机、地址）、重量、费用、状态（已揽收/运输中/已签收）；
  \item 一件快递由一名快递员揽收，一名快递员可以揽收多件快递；
  \item 每个网点有多名快递员，一名快递员只属于一个网点。
\end{itemize}
请完成：
\begin{enumerate}[leftmargin=2em, itemsep=3pt]
  \item[（1）] （5分）画出该系统的 E-R 图。
  \item[（2）] （5分）将 E-R 图转换为关系模式（标明主键和外键）。
\end{enumerate}
\anslines{5cm}

\qtypename{六、应用题}{共2题，每题10分，共20分}

\q 现有 MySQL 数据库中的读者借阅系统部分表：
\begin{lstlisting}[language=SQL]
CREATE TABLE reader (
    rid    INT PRIMARY KEY,
    rname  VARCHAR(50) NOT NULL,
    rtype  VARCHAR(20)  -- 'student' or 'teacher'
);

CREATE TABLE book (
    bid    INT PRIMARY KEY,
    bname  VARCHAR(100),
    bcount INT DEFAULT 1  -- 馆藏数量
);

CREATE TABLE borrow (
    rid     INT,
    bid     INT,
    bdate   DATE,
    rdate   DATE,
    PRIMARY KEY (rid, bid, bdate)
);
\end{lstlisting}
\begin{enumerate}[leftmargin=2em, itemsep=3pt]
  \item[（1）] （3分）查询当前借阅未还（rdate IS NULL）的读者姓名和书名。
  \item[（3）] （3分）查询借阅次数最多的前 5 本书的书名及借阅次数。
  \item[（3）] （4分）创建存储过程 \code{borrow\_book(IN p\_rid INT, IN p\_bid INT)}，实现借书操作：检查该书是否有库存，如果有则插入借书记录（bdate 为当天）并更新库存减 1，否则输出 "库存不足"。要求在过程中使用条件处理和事务控制。
\end{enumerate}
\anslines{4cm}

\q 现有 MySQL 数据库，回答以下问题：
\begin{enumerate}[leftmargin=2em, itemsep=3pt]
  \item[（1）] （3分）简述 MySQL 中 \code{MyISAM} 和 \code{InnoDB} 两种存储引擎的主要区别。
  \item[（3）] （3分）简述 MySQL 的二进制日志（binlog）的作用，它如何帮助数据库恢复？
  \item[（4）] （4分）现有数据表 \code{users(id, username, password)}。编写一条 \code{GRANT} 语句，创建一个新用户 \code{'dbadmin'@'localhost'}，密码为 \code{Admin123!}，并授予其对所有数据库的所有权限。同时写出查看该用户权限的语句和回收所有权限的语句。
\end{enumerate}
\anslines{3.5cm}

\end{document}
